\section{Related Work}\label{sec:related}

This section places our work among five lines of research. For each line it states what others have done, how our
work relates to it and what we reuse. It closes with the gaps that remain open and with Table~\ref{tab:related},
which compares the dependence models that have been proposed or used for quantum programs.

\subsection{Program slicing and its semantics}
Weiser introduced program slicing~\cite{weiser1984program} and showed that programmers use slices when they
debug~\cite{weiser1982programmers}. Slicing extracts a computation that is scattered through a program from the
statements around it, which makes it a comprehension aid as much as a debugging one~\cite{binkley2007source}.
Dependence graphs reduced slicing to graph reachability~\cite{ottenstein1984pdg, ferrante1987pdg,
horwitz1990interprocedural}, Podgurski and Clarke gave dependence a semantic definition~\cite{podgurski1990formal},
and dynamic slicing restricts a slice to one execution~\cite{korel1988dynamic, agrawal1990dynamic}. A backward slice
answers which statements can affect a value at a point, and a forward slice, its dual, answers which statements a
given statement can affect~\cite{tip1994a, horwitz1990interprocedural}. Slicers need not build a graph, and srcSlice,
for example, computes dependences from a slice profile of every variable~\cite{alomari2014srcslice}. Surveys catalog
the many kinds of slices that followed~\cite{tip1994a, silva2012vocabulary, binkley2004a}. Probabilistic slicing
preserves the distribution of the criterion variables~\cite{hur2014slicing, amtoft2020theory, olmedo2025slicing}, and
Gallagher names quantum programs as a direction for semantically defined slicing~\cite{gallagher2025program}.

\emph{Relation and reuse.} We keep the classical notions, a criterion, a backward and a forward slice, and dependence
as the possibility of change, and define them for the reduced state of a set of qubits, a criterion stronger than a
distribution in one basis, which misses phase.

\subsection{Dependence models for quantum programs}
Yu and Zhao's quantum program dependence graph (QPDG)~\cite{yu2025qpdg} extends the classical
PDG~\cite{ferrante1987pdg} with entanglement, no-cloning, temporal and measurement dependences and proposes slicing on
it. It is defined in prose with examples and has not been implemented or evaluated. Its own example shows why
validation matters. In the graph Yu and Zhao draw for their first listing, the slice of a measured qubit omits the two
gates that fix its outcome, so the slice returns 1 with probability \numYzFigSliceProb{} where the program returns 1
with certainty. Kaul et al.\ build the quantum code property graph (QCPG), a code property graph for Qiskit and
OpenQASM programs, and query it for bugs and smells, among them superfluous operations, gates that do not affect a
measured qubit~\cite{kaul2023uniform}. Its data-flow edges run from the control of a CNOT to its target only, and the
authors state that the control is not changed, so the query can report gates whose effect reaches a measured qubit
back through a control. Kaul et al.\ demonstrate the graph on example programs and note the lack of large datasets.
Behler et al.\ derive the order of qubit operations from OpenQASM~3 source and use it to detect patterns and suggest
refactorings, and their order table marks the control of a CNOT apart from its target~\cite{behler2024supporting,
behler2025static}. In their roadmap of quantum software engineering, Murillo et al.\ describe the QPDG as an aid to
understanding programs~\cite{murillo2025roadmap}, and in their roadmap of testing and debugging, Ramalho et al.\ find
dependence and data-flow analysis of circuits almost unexplored, naming only Kaul et al.~\cite{ramalho2025road}.

\emph{Relation and reuse.} We take the QPDG's question, which operations a qubit depends on, and answer it against a
semantic definition. Section~\ref{sec:rules} describes the baseline we reimplement from its prose and the classical
reading of the control that the QCPG states, and Section~\ref{sec:results} measures both.

\subsection{Entanglement analysis}
Static analyses approximate which qubits may be entangled, with type systems~\cite{perdrix2007types,
prost2009reasoning, yuan2022twist}, abstract interpretation~\cite{perdrix2008entanglement, yu2021abstract}, the
stabilizer formalism~\cite{honda2015stabilizer} and data-flow rules~\cite{xia2025entanglement}. ScaffCC assumes that
qubits that interact may be entangled~\cite{javadiabhari2015scaffcc}.

\emph{Relation.} These analyses answer a different question, whether two qubits may be correlated, and they are not
slicers. Section~\ref{sec:core} shows that entanglement neither implies nor is implied by dependence, so their
results cannot decide what a slice keeps.

\subsection{Light cones and compilers}
The \emph{backward light cone} of a set of qubits at the end of a circuit is the set of gates that can affect their
final state. Compilers build gate-dependence graphs on this principle to schedule and map
circuits~\cite{murali2019triq, amy2020staq}, and QAOA analyses use it to bound the gates that can reach an
observable~\cite{farhi2014qaoa}. Quantum separation logics reason locally about the qubits a program
touches~\cite{zhou2021qsl, le2022qsl}. Qiskit implements the light cone as the \texttt{LightCone} transpiler
pass~\cite{qiskit2025}. Given the qubits to be measured, or a Pauli observable on them, the pass walks the circuit
backward from the end, keeps the gates in the light cone and also removes every gate that commutes with the
observable, since such a gate cannot change its expectation value. It returns a smaller circuit for estimating that
observable, decides commutation for the circuit's actual angles and rejects circuits with mid-circuit measurements.
Heim et al.\ likewise argue that quantum dependences should account for commutation~\cite{heim2020quantum}. Compilers
also call the time windows of a circuit its slices~\cite{smith2022timestitch}, a sense different from ours.

\emph{Relation and reuse.} On qubits, the backward light cone is the relation of our write-all rule, and that it
determines the final state of its qubits is known. We claim neither the relation nor its soundness as new. What these
works do not provide is a definition of what a slice must preserve, a proof that the light cone keeps every
dependence, and a measurement of what smaller or different rules lose, which Sections~\ref{sec:core}
and~\ref{sec:results} give. Removing commuting gates goes further than the light cone. A pruned slice is correct for
the observable at the given angles but no longer contains every operation that could change the qubit, so it can omit
the place where a fault would be. We reuse Qiskit's \texttt{LightCone} pass as the commutation-pruning baseline and
follow the light cone along classical bits, which the pass does not do.

\subsection{Debugging tools and bug detectors}
Metwalli and Van Meter adapt slicing to circuits by splitting a circuit at barriers the developer places and removing
the qubits a part does not use~\cite{metwalli2022debugging, metwalli2024testing}, and Qiskit's circuit-slicing
utilities cut a circuit into consecutive parts by barriers, depth or gate type~\cite{qiskitaddonutils}. Both cut a
circuit into pieces to inspect and follow no dependence. Ramalho et al.\ list slicing as a way around no-cloning and
note that it grows hard to manage as circuits grow~\cite{ramalho2025road}. Zappin et al.\ surveyed and interviewed 26
quantum developers, who debug with print statements, circuit drawings, simulators and manual inspection of the code
and report that these methods do not scale~\cite{zappin2026challenges}. Testing faces the limit that phase faults
escape computational-basis measurement~\cite{oldfield2024faster}. Bug detectors find known patterns, and Chen et al.\
catalog quantum-specific code smells~\cite{chen2023smelly}. LintQ runs queries over Qiskit programs and orders two
operations when they share a qubit~\cite{paltenghi2024lintq}, QChecker matches patterns in the abstract syntax
tree~\cite{zhao2023qchecker}, and Q-PAC also detects known patterns~\cite{nayak2023qpac}. Bug
patterns~\cite{zhao2021bugpatterns}, Bugs4Q~\cite{zhao2023bugs4q}, mutation tools~\cite{mendiluze2021muskit,
fortunato2022qmutpy, garcia2022testing} and mutation-based fault localization~\cite{ishimoto2025mbfl} supply faults
and rankings to this field.

\emph{Relation and reuse.} Slices complement these tools. Manual slicing has no criterion, so the developer still
narrows the fault by hand, and a slice by the write-all rule generalizes the removal of unused qubits, since it drops
every qubit and operation that cannot reach the criterion. LintQ's ordering relation is the relation of the write-all
rule, which confirms the rule in practice, while LintQ uses it for patterns and makes no claim about slices. Of the
42 Bugs4Q bugs on which QChecker was evaluated, 22 are wrong outputs that no pattern detects~\cite{zhao2023qchecker},
and a slice shows where such a fault can be. We reuse Bugs4Q, published mutation operators and mutation-based fault
localization.

\subsection{Open gaps and the data used}
Table~\ref{tab:related} compares the models of dependence these lines of work use, by what a CNOT writes, whether
measured bits are followed, which extra edges are added, and what is defined and guaranteed about a slice. Three gaps
remain open, and each can be shown from the work above. First, no dependence rule proposed for quantum slicing has
been checked against what a quantum program computes, and the one example worked out in detail, the QPDG's, gives a
wrong slice. Second, the analyses that represent quantum programs for static analysis, the QCPG and the order tables
of Behler et al., read the control of a gate as unchanged, and no work has measured what that reading does to a slice.
Third, the light cone is known to be sound, but no work has asked whether it keeps every dependence, how far it is from
the smallest slice that does, and what smaller rules lose in correctness and in locating faults.

\begin{table*}[t]
  \centering
  \caption{Dependence models for quantum programs, by what they read a gate to do and what they guarantee}
  \label{tab:related}
  \footnotesize\setlength{\tabcolsep}{3.4pt}\renewcommand{\arraystretch}{1.18}
  \newcommand{\sep}{\arrayrulecolor{black!25}\hline\arrayrulecolor{black}}
  \begin{tabularx}{\textwidth}{@{}>{\raggedright\arraybackslash}p{3.1cm} >{\centering\arraybackslash}p{1.55cm}
      >{\centering\arraybackslash}p{1.25cm} >{\raggedright\arraybackslash}p{2.0cm}
      >{\centering\arraybackslash}p{1.3cm} >{\raggedright\arraybackslash}p{1.8cm} >{\raggedright\arraybackslash}X@{}}
    \toprule
    \textbf{Approach} & \textbf{CNOT writes} & \textbf{Measured bits} & \textbf{Extra edges} & \textbf{Slice semantics} &
    \textbf{Guarantee} & \textbf{Evaluated on} \\
    \midrule
    QPDG, Yu and Zhao~\cite{yu2025qpdg}        & ?        & \full & entanglement, closed & \none & none & 5 listings, not implemented \\ \sep
    QCPG, Kaul et al.~\cite{kaul2023uniform}   & $t$      & \full & none & \none & none & 3 listings, 1 Qiskit/OpenQASM circuit \\ \sep
    Behler et al.~\cite{behler2024supporting, behler2025static} & $t$ & \half & none & \none & none & 16 examples of OpenQASM~3 \\ \sep % TODO verify: QStatic's evaluation set
    LintQ~\cite{paltenghi2024lintq}            & $c, t$   & \full & none & \none & none & \numSizeLintqFiles{} GitHub programs, \numLqTp{} confirmed bugs \\ \sep
    TriQ, staq~\cite{murali2019triq, amy2020staq} & $c, t$ & \na  & none & \none & sound, known & 12 and 24 benchmark circuits \\ \sep
    Qiskit \texttt{LightCone}~\cite{qiskit2025} & $c, t$  & \none & none, removes commuting gates & \half & observable, fixed angles & \na \\
    \midrule
    \rowcolor{blue!7}\textbf{QDG, this paper} & $c, t$ & \full & none & \full & sound, complete & \numNPairs{} circuit--qubit pairs of \numNCircuits{} circuits, \numFamCircuits{} generated and \numSubjLCount{} large circuits, \numRqThreeUnits{} seeded and \numLinesUnits{} Bugs4Q fault units \\
    \bottomrule
  \end{tabularx}\\[3pt]
  \begin{minipage}{\textwidth}\footnotesize
  \full{} has the property, \half{} has it in part, \none{} lacks it, ? means the authors do not specify it, and
  \na{} means it does not apply. \emph{CNOT writes}, the qubits a CNOT with control $c$ and target $t$ is read to
  change, from which the dependences follow. \emph{Measured bits}, classical control is followed from a measurement
  to the operations its bit conditions. \emph{Extra edges}, dependences added beyond those of the gates' qubits.
  \emph{Slice semantics}, what a slice must preserve is defined. \emph{Guarantee}, the property of a slice that the
  authors state or prove; complete means that it keeps every dependence. \emph{Evaluated on}, the data of the original
  evaluation.
  \end{minipage}
\end{table*}

The data differ as much as the models. The dependence graphs were shown on examples and the detectors evaluated on
Bugs4Q and on LintQ's GitHub corpus, while no rule for quantum slicing was evaluated against a ground truth. We use
Bugs4Q for real faults and add what a study of rules needs, circuits with a simulated ground truth and seeded faults,
which Section~\ref{sec:design} lists.
